\section{Introduction}

Auto-research is not just paper generation. A system that produces a plausible PDF has not necessarily conducted a valid research process: it may have proposed a weak question, implemented a different method from the one described, run only a subset of the claimed experiments, copied numbers incorrectly, or cited nonexistent prior work. For end-to-end research agents, the central evaluation problem is therefore not only whether the final manuscript looks convincing, but whether the manuscript is faithful to the research that was actually conducted.

This distinction is increasingly important as general-purpose command-line agents become capable of operating full software environments. CLI agents can read papers, edit code, run experiments, inspect logs, revise drafts, and assemble final PDFs. These abilities make them natural substrates for automated research pipelines, but they also create new failure modes. A final paper can look coherent while its code fails to implement the method, its logs show incomplete runs, or its reported numbers diverge from stored results.

We introduce \ra, a benchmark and study of off-the-shelf CLI agents conducting an end-to-end research pipeline with minimal scaffolding. The benchmark covers three widely used CLI agents--Claude Code, Codex, and Kimi Code--across 13 computer-science domains, including 5 CPU-only and 8 GPU domains. Each agent generates 39 papers, for 117 papers total. For every run, \ra\ records the final paper, code, configurations, execution logs, result artifacts, reviews, and metadata.

We also introduce artifact-grounded peer review. In contrast to PDF-only review, artifact-grounded peer review gives reviewers access to the paper and its process artifacts: code, logs, result files, and references. The goal is not to claim that one automatic reviewer is universally better than another. Rather, PDF-only and artifact-grounded review evaluate different targets. PDF-only review asks whether the final manuscript looks convincing. Artifact-grounded peer review asks whether the paper is faithful to the research process and evidence that produced it.

\paragraph{Contributions.}
First, we present \ra, a large-scale empirical evaluation of off-the-shelf CLI agents conducting end-to-end research-like workflows under a minimal scaffold. The dataset contains 117 generated papers and 351 artifact-grounded peer reviews across 13 CS domains. Second, we define and instantiate artifact-grounded peer review as a process-grounded evaluation protocol for auto-research. The protocol reveals unsupported results, setting mismatches, incomplete experiments, hallucinated references, and weak baselines that are often invisible from the PDF alone.

\paragraph{Summary of findings.}
Current CLI agents can often complete the research loop, but completion does not imply validity. Under Stanford Agentic Reviewer (\sar), a PDF-only reviewer, Claude Code scores 5.45, Codex 4.93, and Kimi Code 4.24 on a 0--10 ICLR-style scale; Claude Code sits near the weighted mean of a sampled ICLR 2025 submission set. Under artifact-grounded peer review (\pr), scores fall to 4.59, 4.51, and 3.38, respectively. The gap is not a contradiction: it reflects a change in evaluation target from apparent paper quality to paper-artifact faithfulness. Across agents, experimental rigor and significance are the dominant bottlenecks. Codex is the strongest on integrity-related dimensions, Claude Code is strongest on novelty and significance, and Kimi Code has the highest mismatch and fake-reference rates. Scaling Codex GPU runs from RTX A6000 to H100 changes \pr\ from about 4.50 to 4.26, suggesting that additional compute alone does not reliably improve auto-research quality.
